% Work/MAT -- 1. rocnik -- hodina 39 -- finalni revize 2026-10-30.
% Priklady z uzivatelem dodanych obrazku upraveny a doplneny vlastnimi ulohami.
% Vysledky a podminky kontrolovany nezavisle; zapis vysledku nepredpoklada
% kladnost promennych tam, kde to neni stanoveno.
\documentclass[11pt,a4paper]{article}
% Spolecny styl pro ucitelske pripravy MAT (Work/MAT)
% Verze 1.1, 2026-10-23
\NeedsTeXFormat{LaTeX2e}
\ProvidesPackage{priprava}[2026/10/23 v1.1 Sablona ucitelskych priprav MAT]

% Zakladni typografie a cestina
\RequirePackage[T1]{fontenc}
\RequirePackage[utf8]{inputenc}
\RequirePackage[czech]{babel}
\RequirePackage{lmodern}
\RequirePackage{microtype}
\RequirePackage{amsmath,amssymb}
\RequirePackage{xcolor}
\RequirePackage{enumitem}
\RequirePackage[most]{tcolorbox}
\RequirePackage{tabularx,booktabs,array}
\RequirePackage{fancyhdr}
\RequirePackage{needspace}
\RequirePackage{geometry}
\RequirePackage[hidelinks]{hyperref}

\geometry{a4paper,top=17mm,bottom=18mm,left=23mm,right=23mm}
\definecolor{notegray}{gray}{0.39}
\setlength{\parindent}{0pt}
\setlength{\parskip}{2.5pt plus 1pt minus 0.5pt}
\setlist[itemize]{topsep=2pt,itemsep=1pt,parsep=0pt,leftmargin=1.45em}
\setlist[enumerate]{topsep=2pt,itemsep=2pt,parsep=0pt,leftmargin=1.65em}
\setcounter{secnumdepth}{0}
\renewcommand{\arraystretch}{1.12}

% Zahlavi: vlevo rocnik | globalni tema, vpravo cislo: nazev hodiny
% Zapati zustava zcela prazdne.
\pagestyle{fancy}
\fancyhf{}
\renewcommand{\headrulewidth}{0.35pt}
\renewcommand{\footrulewidth}{0pt}
\setlength{\headheight}{14pt}
\newcommand{\nastavPripravu}[4]{%
  \fancyhead[L]{\small #1\enspace\textbar\enspace #2}%
  \fancyhead[R]{\small #3: #4}%
}

% Minimalisticke cernobile boxy. Kazda definice zvlast.
\tcbset{commonlesson/.style={enhanced,breakable=false,colback=white,colframe=black,
  colbacktitle=white,coltitle=black,boxrule=0.45pt,
  arc=0.8mm,outer arc=0.8mm,left=3mm,right=3mm,top=1.6mm,bottom=1.6mm,
  toptitle=1mm,bottomtitle=0.7mm,fonttitle=\bfseries,
  attach boxed title to top left={xshift=2.5mm,yshift=-1.15mm},
  boxed title style={colback=white,colframe=white,boxrule=0pt,left=1mm,right=1mm}}}
\newtcolorbox{zadaniBox}[1][]{commonlesson,title={Zadání},#1}
\newtcolorbox{prikladBox}[1][]{commonlesson,title={Příklad},#1}
\newtcolorbox{definiceBox}[1][]{commonlesson,title={Definice},#1}

% Spolecny box pro zadani a jeho reseni, bez mezery mezi castmi.
% Pouziti:
%   \begin{prikladReseniBox}
%     Zadani ulohy.
%     \tcblower
%     \begin{align*} ... \end{align*}
%   \end{prikladReseniBox}
% Spodni cast ma jemny sedy podklad a delici linku; lze v ni sazet
% vice radku matematiky. Obe casti maji jen jeden spolecny ramecek.
\tcbset{solutionlesson/.style={
  commonlesson,
  enhanced jigsaw,
  breakable,
  colbacklower=black!3,
  segmentation style={solid,draw=black!25,line width=0.35pt},
  fontlower=\small\color{notegray},
  before lower={\noindent\textbf{Řešení:}\enspace},
  leftlower=3mm,rightlower=3mm
}}
\newtcolorbox{prikladReseniBox}[1][]{solutionlesson,title={Příklad},#1}
\newtcolorbox{zadaniReseniBox}[1][]{solutionlesson,title={Zadání},#1}

% Pomucky pro autora pripravy: sedive poznamky; \reseni pro starsi pripravy
\newcommand{\poznamka}[1]{\par\begingroup\small\color{notegray}\itshape
  \textnormal{\textbf{Poznámka pro učitele.}}\enspace #1\par\endgroup}
\newcommand{\reseni}[1]{\par\begingroup\small\color{notegray}
  \textbf{Očekávané řešení:}\enspace #1\par\endgroup}
\newcommand{\blok}[1]{\needspace{5\baselineskip}\subsection*{#1}}
\newcommand{\vyrok}[1]{\uv{#1}}

\nastavPripravu{1. ročník}{Mocniny a odmocniny}{39}{Příprava na test a chyby}
\begin{document}

\poznamka{\textbf{Jádro na 45 minut:} 4 min diagnostika znamének a podmínek,
9 min Příklad 1, 12 min Samostatná práce 1 (4 položky),
8 min společná kontrola prvního chybného kroku,
9 min Příklad 2 + Samostatná práce 2 (vyber dvě položky),
3 min shrnutí. Tisknutelný pracovní list je \textbf{samostatný soubor},
ne test s klasifikací. Složitější příklady a domácí úkol jsou nadstavba.}

\blok{Nechybí vzorec; chybu často způsobí špatná strategie}
Před písemkou musejí žáci pracovat samostatně a kontrolovat
\textbf{znaménka, definiční obor, závorky a vhodnost pořadí operací}.
Když výsledek nevychází, požaduj \emph{první nesprávný řádek}, nikoli
pouhou opravu výsledného čísla.

\begin{zadaniReseniBox}[title={Samostatná práce 1: čtyři výroky na úvod},breakable=false]
Urči, zda jsou tvrzení pravdivá pro všechna přípustná čísla; nepravdivá vyvrať.
\begin{enumerate}[label=\alph*)]
\item $\displaystyle\frac{a^3}{a^{-2}}=a^5$ pro $a\ne0$.
\item $\sqrt{(x-2)^2}=x-2$ pro všechna reálná $x$.
\item $\sqrt{18}+\sqrt8=5\sqrt2$.
\item $(a+b)^2=a^2+b^2$ pro všechna reálná $a,b$.
\end{enumerate}
\tcblower
\begin{enumerate}[label=\alph*)]
\item Ano, rozdíl exponentů $3-(-2)=5$.
\item Ne; např. $x=0$ dává vlevo $2$, vpravo $-2$; správně $|x-2|$.
\item Ano: $3\sqrt2+2\sqrt2=5\sqrt2$.
\item Ne; pro $a=b=1$ je $4\ne2$.
\end{enumerate}
\end{zadaniReseniBox}

\blok{1. Najdi první chybu, pak zvol vlastní cestu}
\begin{prikladReseniBox}[title={Příklad 1: tři zavádějící výpočty},breakable=false]
Najdi první chybný krok a uprav výraz správně:
\begin{enumerate}[label=\alph*)]
\item $\displaystyle\frac{(a^{-2})^3a^8}{a^2}=a^{-6+8-2}=a^0=0$.
\item $\displaystyle\frac1{\sqrt7+2}=\frac{\sqrt7+2}{7+4}$.
\item $\sqrt{x^2y^4}=xy^2$ pro všechna reálná $x,y$.
\end{enumerate}
\tcblower
\begin{enumerate}[label=\alph*)]
\item $a^0=\boxed1$, nikoliv $0$; původně $a\ne0$.
\item Vynásob \emph{sdruženým} výrazem:
$\frac{\sqrt7-2}{7-4}=\boxed{(\sqrt7-2)/3}$.
\item $\sqrt{x^2y^4}=\sqrt{(xy^2)^2}=\boxed{|x|y^2}$;
$y^2\ge0$, ale $x$ může být záporné.
\end{enumerate}
\end{prikladReseniBox}

\blok{2. Obtížnější samostatná práce: mocniny a mnoho proměnných}
\begin{zadaniReseniBox}[title={Samostatná práce 2: kombinované algebraické úpravy},breakable=false]
Uprav, nevyčísluj zbytečně velké mocniny. Připiš podmínky.
\begin{enumerate}[label=\alph*)]
\item $\displaystyle\frac{a^{-3}b^2}{c^{-3}}:
\left(\frac{b^3c}{a^2}\right)^2$.
\item $\displaystyle\left(\frac{9a^{-2}c^3}{8b^{-2}}\right)^{-3}$.
\item $\displaystyle\frac{3(a^3b^2)^3}{(2a^2b)^2}\cdot
\left(\frac{2a^4b^3}{a^3b}\right)^3$.
\item $\displaystyle\frac{(x+y)^3(x+y)^4}{(x+y)^5}:(x+y)^2$.
\end{enumerate}
\tcblower
\begin{enumerate}[label=\alph*)]
\item $\boxed{ac/b^4}$, $a,b,c\ne0$.
\item $\boxed{512a^6/(729b^6c^9)}$, $a,b,c\ne0$.
\item $\boxed{6a^8b^{10}}$, $a,b\ne0$.
\item $\boxed1$, $x+y\ne0$.
\end{enumerate}
\end{zadaniReseniBox}

\poznamka{V povinné části stačí dvě položky Samostatné práce 2.
Další jsou diferenciace; žáci s rychlým řešením ať místo kalkulačky
provedou zpětnou kontrolu dosazením například $a=2,b=3,c=5$.}

\blok{3. Odmocniny, sdružené výrazy a kontrola}
\begin{prikladReseniBox}[title={Příklad 2: proč zvolit tento pořad úprav?},breakable=false]
Bez kalkulačky zjednoduš
\[
\frac{\sqrt{75}-\sqrt{12}}{\sqrt3},
\qquad
\frac1{2-\sqrt3}-\frac1{2+\sqrt3}.
\]
\tcblower
První výraz nejprve částečně odmocni:
$\frac{5\sqrt3-2\sqrt3}{\sqrt3}=\boxed3$.
U druhého usměrni oba zlomky (oba jmenovatele mají součin $1$):
$(2+\sqrt3)-(2-\sqrt3)=\boxed{2\sqrt3}$.
\end{prikladReseniBox}

\begin{zadaniReseniBox}[title={Samostatná práce 3: pět úloh bez nápovědy},breakable=false]
\begin{enumerate}[label=\alph*)]
\item $\sqrt{200}-2\sqrt{18}+\sqrt8$.
\item $\sqrt[3]{-250}+2\sqrt[3]{16}$.
\item $\displaystyle\frac3{\sqrt6+\sqrt3}$.
\item $81^{-3/4}\cdot27^{2/3}$.
\item Pro která $x\in\mathbb R$ je definován výraz
$\displaystyle\frac1{\sqrt{x-4}}$?
\end{enumerate}
\tcblower
\begin{enumerate}[label=\alph*)]
\item $\boxed{6\sqrt2}$.
\item $\boxed{-\sqrt[3]2}$.
\item $\boxed{\sqrt6-\sqrt3}$, protože
$(\sqrt6+\sqrt3)(\sqrt6-\sqrt3)=3$.
\item $\boxed{1/3}$.
\item $\boxed{x>4}$; argument odmocniny musí být \emph{kladný},
protože je celé odmocnítko ve jmenovateli.
\end{enumerate}
\end{zadaniReseniBox}

\blok{4. Složitější rezerva: matematika bez chytrého kalkulování}
\begin{zadaniReseniBox}[title={Rozšiřující příklad 1: rozklad na prvočísla},breakable=false]
\begin{enumerate}[label=\alph*)]
\item $\displaystyle\frac{125^2\cdot2^{10}}{16^2\cdot50}\cdot\frac{25}{5^3}$.
\item $\displaystyle\left(\frac{2^2}{5}\right)^5
\left(-\frac{5^2}{2^3}\right)^3$.
\item $\displaystyle\frac{8^2\cdot6^3}{4^3\cdot12}$.
\end{enumerate}
\tcblower
\begin{enumerate}[label=\alph*)]
\item $\boxed{250}$.
\item $\boxed{-10}$.
\item $\boxed{18}$.
\end{enumerate}
\end{zadaniReseniBox}

\blok{Domácí úkol: samostatná kontrola před dvacetiminutovým testem}
\begin{zadaniReseniBox}[title={Domácí úkol 1: osmička na závěrečné procvičení},breakable=false]
\begin{enumerate}[label=\alph*)]
\item $\displaystyle\frac{(-3)^6(-3)^{-2}}{9}$.
\item $\displaystyle\frac{a^{-4}b^5}{a^2b^{-3}}$.
\item $\displaystyle\frac{(2a^3b^2)^3}{4a^4b}$.
\item $\sqrt{288}-2\sqrt{32}+\sqrt{18}$.
\item $\displaystyle\frac1{\sqrt5-\sqrt3}$.
\item $\sqrt[3]{-1024}+\sqrt[3]{128}$.
\item $0{,}000000084\,\mathrm m$ zapiš vědecky a v nanometrech.
\item Soubor má $3{,}5\,\mathrm{MiB}$. Kolik je to $\mathrm{KiB}$?
\end{enumerate}
\tcblower
\begin{enumerate}[label=\alph*)]
\item $(-3)^4/9=81/9=\boxed9$.
\item $\boxed{b^8/a^6}$, $a,b\ne0$.
\item $\boxed{2a^5b^5}$, $a,b\ne0$ (původní jmenovatel vyžaduje oba nenulové).
\item $(12-8+3)\sqrt2=\boxed{7\sqrt2}$.
\item $\boxed{(\sqrt5+\sqrt3)/2}$.
\item $-8\sqrt[3]2+4\sqrt[3]2=\boxed{-4\sqrt[3]2}$.
\item $\boxed{8{,}4\cdot10^{-8}\,\mathrm m=84\,\mathrm{nm}}$.
\item $3{,}5\cdot1024=\boxed{3584\,\mathrm{KiB}}$.
\end{enumerate}
\end{zadaniReseniBox}

\poznamka{Samostatně tisknutelný list \texttt{39\_MOC\_PracovniList\_Procvicovani.pdf}
obsahuje další úlohy a \textbf{výsledky až za příkazem \texttt{\\newpage}}.
Žákům tedy vytiskni pouze stránky před oddílem Výsledky; celý soubor
si nech jako učitelský klíč.}
\end{document}
